% =============================================================================
%  Fully developed sPTT pipe flow with a Newtonian solvent: analytical
%  solution used to verify the solver.  For \input in the article (appendix or
%  verification subsection).  Citation keys: OliveiraPinho1999,
%  Cruz2005, Pacheco2023 (see sptt_pipe_solution.bib).
%  Notation: z axial, r radial, R = D/2, polymer extra stress \boldsymbol{\sigma}.
% =============================================================================

\subsection{Fully developed pipe flow of the sPTT fluid}\label{sec:sptt-pipe}

The solver is verified against the steady, fully developed flow of the
constitutive model of \S\,2 in a straight pipe of radius $R=D/2$. The total
extra stress is $\boldsymbol{T}=2\eta_s\boldsymbol{D}+\boldsymbol{\sigma}$,
with the polymer contribution obeying the linear simplified
Phan-Thien--Tanner law
\begin{equation}
  f(\operatorname{tr}\boldsymbol{\sigma})\,\boldsymbol{\sigma}
  + \lambda\,\overset{\triangledown}{\boldsymbol{\sigma}} = 2\eta_p\boldsymbol{D},
  \qquad
  f = 1 + \frac{\varepsilon\lambda}{\eta_p}\operatorname{tr}\boldsymbol{\sigma},
  \label{eq:sptt}
\end{equation}
where $\overset{\triangledown}{\boldsymbol{\sigma}}$ is the upper-convected
derivative; $\varepsilon=0$ recovers the Oldroyd-B model.

For $\boldsymbol{u}=u(r)\,\boldsymbol{e}_z$ the only non-zero velocity
gradient is the shear rate $\dot\gamma=\mathrm{d}u/\mathrm{d}r\le 0$, the
material derivative of $\boldsymbol{\sigma}$ vanishes and
\eqref{eq:sptt} reduces to
\begin{equation}
  f\,\sigma_{rz} = \eta_p\dot\gamma, \qquad
  f\,\sigma_{zz} = 2\lambda\dot\gamma\,\sigma_{rz}, \qquad
  \sigma_{rr}=\sigma_{\theta\theta}=0 ,
  \label{eq:sptt-components}
\end{equation}
so that
\begin{equation}
  \sigma_{zz} = \frac{2\lambda}{\eta_p}\,\sigma_{rz}^2, \qquad
  f = 1 + a\,\sigma_{rz}^2, \qquad a = \frac{2\varepsilon\lambda^2}{\eta_p^2}.
  \label{eq:sptt-normal}
\end{equation}
The axial momentum balance with pressure gradient $\mathrm{d}p/\mathrm{d}z=-G$
fixes the total shear stress,
\begin{equation}
  \eta_s\dot\gamma + \sigma_{rz} = \tau(r) = -\tfrac12 G r .
  \label{eq:momentum}
\end{equation}
Eliminating $\dot\gamma=f\sigma_{rz}/\eta_p$ between
\eqref{eq:sptt-components} and \eqref{eq:momentum} gives, with
$\kappa=\eta_s/\eta_p$, a cubic for the polymer shear stress
$s=\sigma_{rz}$ \citep{Cruz2005},
\begin{equation}
  \kappa a\, s^3 + (1+\kappa)\, s - \tau = 0 .
  \label{eq:cubic}
\end{equation}
Its left-hand side is strictly increasing in $s$, so \eqref{eq:cubic} has a
single real root, given by Cardano's formula
\begin{equation}
  s = \sqrt[3]{-\frac{q}{2}+\sqrt{\Delta}} + \sqrt[3]{-\frac{q}{2}-\sqrt{\Delta}},
  \label{eq:cardano}
\end{equation}
with
$p = (1+\kappa)/(\kappa a)$, $q = -\tau/(\kappa a)$ and
$\Delta = q^2/4 + p^3/27 > 0$.
The shear rate is $\dot\gamma = (1+a s^2)\,s/\eta_p$, and velocity and flow rate
follow from single quadratures,
\begin{equation}
  u(r) = -\int_r^R \dot\gamma(r')\,\mathrm{d}r', \qquad
  Q = 2\pi\int_0^R u\,r\,\mathrm{d}r = -\pi\int_0^R r^2\dot\gamma\,\mathrm{d}r ,
  \label{eq:velocity}
\end{equation}
the second by parts. Since $Q(G)$ is monotone, the pressure gradient for a
prescribed flow rate is obtained by bisection; we evaluate
\eqref{eq:velocity} with 64-point Gauss--Legendre quadrature, which reproduces
the closed-form limits below to machine precision. The wall shear stress is
$\tau_w = GR/2$ for any rheology; the rheology enters through $G(Q)$.

Two limits have closed forms. For $\varepsilon=0$ (Oldroyd-B) the velocity is the
Poiseuille profile of the total viscosity $\eta_0=\eta_s+\eta_p$,
$u = 2\bar u\,(1-r^2/R^2)$ with $G=32\eta_0\bar u/D^2$, and
$\sigma_{rz}=\eta_p\dot\gamma$, $\sigma_{zz}=2\lambda\eta_p\dot\gamma^2$.
Without solvent ($\eta_s=0$), $s=\tau$ and \citep{OliveiraPinho1999}
\begin{equation}
  u(r) = \frac{G}{4\eta_p}\,(R^2-r^2)
       + \frac{\varepsilon\lambda^2G^3}{16\,\eta_p^3}\,(R^4-r^4),
  \qquad
  Q = \frac{\pi G R^4}{8\eta_p} + \frac{\pi\varepsilon\lambda^2 G^3 R^6}{24\,\eta_p^3}.
  \label{eq:op}
\end{equation}

\paragraph{Test set-up.}
A straight pipe of length $14D$ (inlet at $z=-2D$, outlet at $z=12D$) is
discretised with quasi-uniform tetrahedral meshes of size
$h/D = 0.143,\ 0.1,\ 0.071,\ 0.05$. The fluid has $\beta=\eta_s/\eta_0=0.5$,
$\mathit{Wi}=\lambda\bar u/D=0.5$ and $\varepsilon\in\{0,\,0.25\}$, at
$\mathit{Re}=10$. The inlet imposes the Newtonian parabola and a relaxed
polymer ($\boldsymbol{\psi}=\boldsymbol{0}$), so the flow must develop its
viscoelastic profile inside the domain; the outlet carries a zero traction.
The flow rate is ramped up over the first of two periods
$T=\pi\rho D^2/(2\eta_0)$ ($\mathit{Wo}=1$), with 20 BDF1 steps per period;
the final time is $3.1\,\rho D^2/\eta_0 = 63\lambda$, and the Newton increment
of $\boldsymbol{\psi}$ at the last step is at round-off level, so the
solution compared is the steady one. Relative $L^2$ errors of $\boldsymbol{u}$
and of the full tensor $\boldsymbol{\sigma}$ are integrated over the developed
region $5\le z/D\le 10$, where the reference solution is evaluated at the flow
rate carried by the discrete solution; the pressure gradient $G_h$ is the
least-squares slope of $p_h$ over the same region, and $\tau_{w,h}$ is the
wall shear stress at the wall nodes of the region. Table~\ref{tab:sptt-pipe} lists the errors and the
observed orders $p=\ln(e_{k-1}/e_k)/\ln(h_{k-1}/h_k)$.

% Written by run_sptt_pipe.sh as $OUT/table.tex: replace this file with it.
% PRELIMINARY (container, 2 ranks): the two coarsest meshes only.
\begin{table}
  \centering
  \begin{tabular}{cccccccc}
    \toprule
    $h/D$ & $e_{\boldsymbol{u}}$ & $p$ & $e_{\boldsymbol{\sigma}}$ & $p$
      & $|G_h/G-1|$ & $p$ & $\tau_{w,h}/\tau_w-1$ (\%) \\
    \midrule
    \multicolumn{8}{l}{Oldroyd-B ($\varepsilon=0$)} \\
    0.200 & $6.58{\times}10^{-2}$ & -- & $4.62{\times}10^{-1}$ & -- & $1.65{\times}10^{-1}$ & -- & $+10.9$ \\
    0.143 & $3.58{\times}10^{-2}$ & 1.81 & $1.73{\times}10^{-1}$ & 2.93 & $7.56{\times}10^{-2}$ & 2.31 & $+0.8$ \\
    \multicolumn{8}{l}{sPTT ($\varepsilon=0.25$)} \\
    0.200 & $6.38{\times}10^{-2}$ & -- & $1.39{\times}10^{-1}$ & -- & $1.15{\times}10^{-1}$ & -- & $-4.4$ \\
    0.143 & $3.11{\times}10^{-2}$ & 2.13 & $7.37{\times}10^{-2}$ & 1.88 & $5.44{\times}10^{-2}$ & 2.22 & $-4.3$ \\
    \bottomrule
  \end{tabular}
  \caption{Fully developed pipe flow: relative $L^2$ errors of velocity and
  polymer stress over $5\le z/D\le10$, error of the pressure gradient and mean
  error of the wall shear stress, with observed orders $p$.}
  \label{tab:sptt-pipe}
\end{table}


\paragraph{Remark on manufactured solutions.}
The closed-form unsteady family of \citet{Pacheco2023} for the Oldroyd-B
model, with $\sigma_{11}=\sigma_{22}=\sigma(t)$, $\sigma_{12}=0$ and the
time-dependent relaxation time $\lambda(t)=(\beta-1)\nu_0/\sigma(t)$, cannot be
used with the log-conformation formulation: it implies
$\lambda\boldsymbol{\sigma}/\eta_p=-\boldsymbol{I}$, that is, a conformation
tensor $\boldsymbol{c}=\boldsymbol{I}+\lambda\boldsymbol{\sigma}/\eta_p\equiv\boldsymbol{0}$,
which has no logarithm. Their remaining test, the ramped-up Oldroyd-B channel
flow compared with its steady state (their \S\,5.1.4), is the planar
counterpart of the $\varepsilon=0$ case above.
